%% Unofficial SPI cover LaTeX template
%% 2017-12-07
%
% By Quentin Ayoul-Guilmard
% Contact: quentin.ayoul-guilmard@centraliens-nantes.net
%
%% DISCLAIMER
%
% This template aims at reproducing the look of SPI doctorate school's Microsoft Word template, as it is today (see date above).
% There are minute and purposeful discrepancies, which I consider to be improvements.
% I will not tell where.
%
% To use this template, follow these steps
% 1_ Copy its preamble into yours (except for document class, obviously)
% 2_ Copy the logo files into your document's folder.
% 3_ Copy the content between "FRONT COVER" and "END OF FRONT COVER" at the beginning of your document
% 4_ Copy the content between "BACK COVER" and "END OF BACK COVER" at the end of your document
% 5_ Find all the comments labelled "TODO" and follow their instructions (if relevant).
% Alternatively to steps 3 and 4, you may put the contents in two files front-cover.tex and back-cover.tex, then load them with \input.
%
% Notice: This document has been successfully compiled with pdflatex and TexLive 2016.
% I cannot guarantee that this code will work with your document.
% Furthermore, it will have to be updated as SPI changes its cover.
% I strove to make it easy to understand and modify with my comments.
% Any feedback or improvement is welcome.
% See below for some ideas.
%
% Ideas for improvement:
% _ Change language locally with babel for English abstract.
% _ Several spacing macros have to be adjusted manually (see below). It would be nice to have the corresponding spaces automatically adjusted.
% _ Two other packages (that I know of) can be used to set background pictures: background and wallpaper. Are they more suited than eso-pic here?
%
%% END OF DISCLAIMER

\documentclass[10pt]{article}

% Standard encoding and language packages.
\usepackage[utf8]{inputenc}
\usepackage[T1]{fontenc}
\usepackage[francais]{babel}
% TODO: replace the above packages with your own encoding settings.
% If you have none, you may want to consider using these.

% Required for cover
\usepackage[a4paper]{geometry} % to change margins on cover pages
\usepackage[scaled]{helvet} % font used on cover (Helvetica)
\usepackage{eso-pic} % to set background picture
\usepackage{multicol} % for back cover (abstracts)
\usepackage{graphicx} % to include logos
\usepackage{tikz} % to compose background picture

% Used for template demonstration only
\usepackage{lipsum} % to generate random text with \lipsum.
%TODO: You can remove this package if you use no \lipsum command.

% Colors (extracted from SPI's template)
\definecolor{boxcolor1}{rgb}{0.91373,0.92941,0.87451}
\definecolor{boxcolor2}{rgb}{0.94902,0.93333,0.91373}
\definecolor{boxcolor3}{rgb}{0.76078,0.87843,0.17647}
\definecolor{headercolor}{rgb}{0.94118,0.30980,0.17255}
\definecolor{namecolor}{rgb}{1.0,0.4,0.0}
\definecolor{titlecolor}{rgb}{0.19216,0.51765,0.60784}
% Also used: gray, teal (predefined by xcolor package, usually loaded by document class)

% Cover environment, to keep changes local
\newenvironment{cover}{%
  \fontfamily{phv}\selectfont % Select Helvetica font
  \pagestyle{empty} % No page number
}{%
  % Do not count this page (useless on back cover, obviously)
  \addtocounter{page}{-1}
  % Go to next odd page (i.e. recto, or right page)
  \cleardoublepage
  % This last command is compulsory even for back cover;
  % without it, the page is not completed and some page-wide
  % commands fail to apply.
}

% Macro for background common to front and back
\newcommand{\tikzBG}{%
  \path (0,0) rectangle (1,1);
  %TODO: You should adjust the bottom height of the following rectangle to fit your abstract's length
  \path [fill=boxcolor1] (.0571,.11) rectangle (.481,.963); 
  \path [fill=boxcolor2] (.4333,.697) rectangle (.9048,.7475);
  \path [fill=boxcolor2] (.4333,.7811) rectangle (.9048,.8316);
  \path [fill=boxcolor2] (.4333,.8687) rectangle (.9048,.9192);
  \path [fill=boxcolor3] (.0571,.7879) rectangle (.5762,.8316);
  \node[inner sep=0pt] at (0.2285,0.8788) [above left] {%
    \includegraphics[height=.0707\paperheight,keepaspectratio]{logo_unb}};
  \node[inner sep=0pt] at (0.6667,0.8788) [above right] {%
    \includegraphics[height=.0808\paperheight,keepaspectratio]{logo_ecn_color}};
  \node at (.0571,.8316) [above right,color=headercolor] {%
    \fontsize{29}{35}\selectfont\bfseries Th\`ese de Doctorat};
}

% Macro for repeated information (to avoid insconsistency)
%TODO: fill in with no formatting but desired case
\newcommand{\firstName}{First Name}
\newcommand{\surname}{Family Name}
\newcommand{\thesisTitle}{Thesis title}

\begin{document}

%% FRONT COVER
\begin{cover}
  % Background
  \AddToShipoutPictureBG*{% from eso-pic package
    \begin{tikzpicture}[x=\paperwidth,y=\paperheight]
      % Background common to front and back:
      \tikzBG{}
      % Only on front cover:
      \node at (.66905,.72225) [color=namecolor] {%
        \fontsize{24}{29}\selectfont \firstName{} \textsc{\surname}};
      % TODO: you may want to change font size above according to text length.
    \end{tikzpicture}%
    % Should SPI change its cover's background, you may use an \inclugraphics
    % to include it (as a temporary fix), instead of the TikZ macros in \tikzBG.
  }

  % Front cover margins
  % TODO: If you modified the main box's size, you may have to change
  % the bottom margin for the jury table to be suitably aligned.
  \newgeometry{left=1.5cm, right=2cm, top=9.5cm, bottom=3.5cm}
  
  \noindent
  \begin{minipage}{8.4cm} % in cm: width = 9.9 - left margin
    \noindent
    Mémoire présenté en vue de l'obtention du grade de \mbox{\emph{Docteur de l'\'Ecole Centrale de Nantes}} sous le sceau de l'Université Bretagne Loire
    % \mbox is used to prevent hyphenation
    \bigskip

    \noindent
    \textbf{\'Ecole doctorale} : Sciences pour l'Ingénieur (SPI)
    \medskip

    \noindent
    \textbf{Discipline} : %TODO: name your discipline here.
    \medskip

    \noindent
    \textbf{Unité de recherche} : %TODO: name your laboratory here.
    \medskip

    \noindent
    Soutenue le %TODO: put your date of defense here.
  \end{minipage}

  % Elastic spacing with weight 1
  \vspace{\stretch{1}}

  \begin{center}
    \scshape\color{titlecolor}%
    {\LARGE\bfseries\thesisTitle\bigskip}
    % TODO: you may want to change font size above (and below) according to text length.

    {\large Thesis subtitle}
    % TODO: Replace text above by your subtitle.
    % If you have none, remove it, and the \bigskip too.
  \end{center}

  % Elastic spacing with weight 2
  \vspace{\stretch{2}}

  %TODO: replace with your jury.
  \noindent
  \begin{tabular}{llll} % use {@{}llll@{}} to remove left and right spaces 
    \multicolumn{4}{c}{\scshape\color{teal}\Large Jury} \\[10pt]
    Rapporteurs & Name \textsc{Surname} & Rank & Affiliation\\
                & Name \textsc{Surname} & Rank & Affiliation \\[10pt]
    Examinateurs & Name \textsc{Surname} & Rank & Affiliation
    \\
                & Name \textsc{Surname} & Rank & Affiliation \\
                & Name \textsc{Surname} & Rank & Affiliation \\[10pt]
    Directeur de thèse & Name \textsc{Surname} & Rank & Affiliation \\
    Co-directeur de thèse & Name \textsc{Surname} & Rank & Affiliation \\
    Co-encadrant de thèse & Name \textsc{Surname} & Rank & Affiliation 
  \end{tabular}
\end{cover}
%% END OF FRONT COVER


% TODO: Replace this with your own content. This is for demonstration only.
\section*{Your content with normal margins and font}
\subsection*{Lorem ipsum}
\lipsum
\pagebreak


%% BACK COVER
\begin{cover}

  % TODO: use macro to break to even page (i.e. verso, or left page).
  % This depends on your document class. Some examples below.
  % \cleardoubleevenemptypage % KOMA-script classes
  % \cleartoverso % memoir class
  % If you had to use another one, please contribute this information (see contact info at beginning of file)

  \AddToShipoutPictureBG*{%
    \begin{tikzpicture}[x=\paperwidth,y=\paperheight]
      % Background common to front and back:
      \tikzBG{}
      % Only on back cover:
      \path [fill=boxcolor3] (.49048,.037037) rectangle (.9048,.08081) node at (.69764,.05892) {\scshape\Large\color{white} Université Bretagne Loire};
    \end{tikzpicture}%
  }

  % Back cover margins (different from front cover)
  \newgeometry{left=1.4cm, right=2.1cm, top=6.6cm, bottom=2.6cm}

  % Name
  \noindent\textcolor{gray}{\LARGE \firstName{} \textsc{\surname}}
  % TODO: you may want to change font size above according to text length.
  \vspace{18pt}
  % TODO: You may need to reduce the value above if your French title (see below)
  % spans two lines. Indeed, the French title should (ideally) be vertically centered
  % on the gray box.

  % Titles
  \begin{center}
    \begin{minipage}[c]{.75\linewidth}
      \Large\color{gray}\bfseries
      % TODO: you may want to change font size above according to text length.
      \thesisTitle
      \vspace{18pt}
      % TODO: You may want to reduce the value above if your French title
      % spans two lines. This value was set so high to prevent the English
      % title from partially overlapping the gray box.

      \noindent Thesis title (translated)
      %TODO: Replace text above by your thesis title in English.
    \end{minipage}
  \end{center}
  \vspace{20pt}

  \begin{multicols}{2}
    \noindent\textbf{\large Résumé}
    \medskip

    \small\noindent % no need to indent first paragraph
    %TODO: Replace the placeholder below by your French abstract.
    \lipsum[1-3]
    \bigskip

    \noindent
    \textbf{Mots-clefs}
    \smallskip

    \noindent %TODO: put your French keywords here.
    % Suggested format:
    keyword1 ; keyword2 ; keyword3 ; \dots{} ; final keyword.

    % Go to next column
    \columnbreak

    \noindent
    \textbf{\large Abstract}
    \medskip
    
    \noindent % no need to indent first paragraph
    %TODO: Replace the placeholder below by your English abstract.
    \lipsum[4-6]
    \bigskip

    \noindent
    \textbf{Keywords}
    \smallskip

    \noindent %TODO: put your English keywords here.
    % Suggested format:
    keyword1; keyword2; keyword3; \dots{}; final keyword.
    
  \end{multicols}
\end{cover}

\end{document}

%%% Local Variables:
%%% mode: latex
%%% TeX-master: t
%%% End:
